\documentclass[11pt]{article}

\usepackage[letterpaper,margin=1in]{geometry}
\usepackage{microtype}
\usepackage{amsmath}
\usepackage[numbers,sort&compress]{natbib}
\usepackage{booktabs}
\usepackage{array}
\usepackage{float}
\usepackage{graphicx}
\usepackage{xcolor}
\usepackage[most]{tcolorbox}
\usepackage{enumitem}
\usepackage{caption}
\usepackage[hidelinks]{hyperref}
\usepackage{parskip}
\usepackage[nameinlink,capitalize]{cleveref}
\setlength{\emergencystretch}{3em}
\graphicspath{{figures/}}

% ---- teaching / takeaway callout boxes (reused from the simplified editions) ----
\definecolor{bgblue}{HTML}{EAF2FF}
\definecolor{bgblueedge}{HTML}{2563EB}
\definecolor{tkgreen}{HTML}{EAF7EE}
\definecolor{tkgreenedge}{HTML}{15803D}
\newtcolorbox{background}[1]{colback=bgblue,colframe=bgblueedge,title={#1},fonttitle=\bfseries\small,
  fontupper=\small,boxrule=0.5pt,left=4pt,right=4pt,top=3pt,bottom=3pt,breakable}
\newtcolorbox{takeaway}[1]{colback=tkgreen,colframe=tkgreenedge,title={#1},fonttitle=\bfseries\small,
  fontupper=\small,boxrule=0.5pt,left=4pt,right=4pt,top=3pt,bottom=3pt,breakable}
\newcommand{\code}[1]{\texttt{\small #1}}

% ---- study-status banner (reused from Paper B) ----
\definecolor{statusred}{HTML}{9F1239}
\definecolor{statusfill}{HTML}{FFF1F2}
\newcommand{\draftwarning}[1]{%
  \begin{center}
  \fcolorbox{statusred}{statusfill}{%
    \begin{minipage}{0.95\linewidth}\small\textbf{Study status.} #1\end{minipage}}
  \end{center}}

% ---- pending-result marker (for the objective axis + ExpGuard, until their locked runs land) ----
\newtcolorbox{pending}[1]{colback=black!4,colframe=black!45,title={#1},fonttitle=\bfseries\small,
  fontupper=\small,boxrule=0.5pt,left=4pt,right=4pt,top=3pt,bottom=3pt,breakable}

% ---- generated macros bound to the locked results (numbers are never hand-typed) ----
\input{generated/results_macros_gen} % Paper A: \Rep*, \Transfer*, \HarmBench*, \SpecializationSeedCount, ...
\input{generated/pilot_macros}       % Paper B: \Pilot*

\title{\textbf{The Benchmark Chooses the Winner}\\[3pt]
\large Honestly Measuring, Tuning, and Composing Small Safety Guards\\
Across Objectives and Four High-Compliance Domains (Law $\cdot$ Finance $\cdot$ Health $\cdot$ Mortgage)}
\author{Reza Rahimi\\\small JazzX AI, Los Altos, CA, USA \\ \small\texttt{reza.rahimi@jazzx.ai}}
\date{Working draft --- \today}

\begin{document}
\maketitle

\draftwarning{This report unifies three studies and two new experiment tracks, and it is written to
be read honestly. \textbf{Two evidence flavors, never pooled:} (1) the fine-tuning / composition /
objective-axis results are \emph{retrospective, estimation-only} on a fixed 4-checkpoint panel and a
manifest inspected during development --- \emph{not} confirmatory; (2) the mortgage benchmark's labels
are \emph{LLM-judge, policy-card-consistent --- not SME-adjudicated}, while the finance/health/law
numbers come from the external, \emph{expert-annotated} ExpGuard set. No causal, universal, or
fair-lending claim is licensed. Sections marked \textsc{pending} await their locked run and contain no
invented numbers.}

\tableofcontents

\begin{abstract}
Teams routinely turn a small chat model into a binary safety \emph{guard} with a short fine-tune and
report its score on a public benchmark. We argue that score is not a property of the guard: \emph{the
benchmark co-produces the verdict.} We show this at three scales on one fixed panel of four compact
checkpoints (Qwen2.5-1.5B, SmolLM2-1.7B, SmolLM3-3B, Qwen3-4B). \textbf{(I)} Supervised fine-tuning
(SFT) raises represented-source ranking by $\RepDelta$ macro-AP [$\RepCILower,\RepCIUpper$] but changes
dataset-held-out transfer by only $\TransferDelta$ [$\TransferCILower,\TransferCIUpper$]: the guard
\emph{specializes}. \textbf{(II, \textsc{pending})} We ask whether the fine-tuning \emph{objective}
(SFT vs DPO vs GRPO) changes that trade-off, pre-registering GRPO's likely single-token null.
\textbf{(III)} A retraining-free \emph{composition} --- averaging the base and fine-tuned calibrated
scores --- recovers transfer relative to SFT ($\PilotTransferDeltaSFT$
[$\PilotTransferDeltaSFTLow,\PilotTransferDeltaSFTHigh$]) at a small represented cost
($\PilotRepDeltaSFT$). \textbf{(IV)} Across four high-compliance domains, a general-safety score does
not cover domain compliance: our HMDA-grounded, dual-labeled \emph{mortgage} benchmark isolates
requests that read as safe yet violate policy (the G0/D1 stratum), and an external replication on
\emph{finance/health/law} (ExpGuard) tests whether the specialization pattern recurs. We ship a single
\code{reproduce} entry point that regenerates every number here from committed per-row scores.
\end{abstract}

\section{Introduction: why ``is this guard good?'' is the wrong question}
\label{sec:intro}

\begin{background}{In one paragraph}
A ``guard'' is a small model that reads an incoming request and labels it \code{safe} or \code{unsafe}
before an assistant acts on it. The common recipe is: take a general chat model, fine-tune it into a
guard, and report its benchmark score --- usually compared against \emph{other} models. That hides the
quantity a practitioner actually needs: what did the fine-tune change relative to the \emph{same model
before tuning}, and does that change survive on data the guard never trained on? This report answers
that at three scales (does tuning transfer? does the \emph{objective} matter? can we compose to
recover?) and then asks whether any of it holds where it matters most --- \emph{regulated domains}
where a violation can read as perfectly polite text.
\end{background}

Converting a general instruction model into a safety guard with a short parameter-efficient fine-tune
\citep{elesedy2024loraguard} and reporting its score on a public suite \citep{bassani2024guardbench}
is now standard. A post-tuning leaderboard number, however, conflates two different things: raising the
score on benchmark \emph{sources represented in training} versus improving \emph{transfer} to datasets
the guard never saw. This report is organized around one thesis --- \emph{the benchmark co-produces the
verdict} --- and four questions on a single shared panel of checkpoints:
\begin{enumerate}[leftmargin=1.6em,itemsep=1pt]
  \item \textbf{Does fine-tuning transfer?} (\Cref{sec:actI}) --- or does the guard specialize to its training sources?
  \item \textbf{Does the objective matter?} (\Cref{sec:actII}) --- SFT vs DPO vs GRPO. \textsc{pending} the locked run.
  \item \textbf{Can we recover transfer without retraining?} (\Cref{sec:actIII}) --- by composing base $+$ adapter.
  \item \textbf{Does any of it survive a change of domain?} (\Cref{sec:actIV}) --- across law, finance, health, and mortgage.
\end{enumerate}

\paragraph{What is new here (and what is not).} Prior work shows fine-tuning can weaken safety behavior
\citep{hsiung2025guardrails}, that guards overfit their training policy
\citep{liu2026domaingeneralizable}, and that a plain instruction model can rival a purpose-built guard
\citep{bassani2024guardbench}; those compare \emph{different} models. Our through-line is a
\emph{paired, same-checkpoint} measurement --- each guard against its own base, on identical rows and
scoring --- extended across \emph{objectives}, a \emph{composition} remedy, and \emph{four domains}.
We introduce no new model, metric, or training algorithm, and make no causal or universal claim; the
output is a reproducible, estimation-only characterization of this fixed panel plus one external
domain replication.

\paragraph{Scope.} English input-prompt classification; binary \code{safe}/\code{unsafe}; four compact
checkpoints (1.5--4B); one LoRA recipe; objectives SFT/DPO/GRPO; four regulated domains. Not a
leaderboard, not a fair-lending audit, not a deployment recommendation.

\section{Background you'll need}
\label{sec:background}

\begin{background}{Terms, once, with a worked example}
\textbf{Guard \& score.} The guard reads a prompt and we look at the single next word it would emit,
comparing its score for \code{unsafe} vs \code{safe}: $s(x)=z_{\text{unsafe}}-z_{\text{safe}}$ (a
``logit'' is the model's raw pre-probability score). Positive $\Rightarrow$ leans unsafe. A two-way
softmax turns $s$ into a probability (e.g.\ $z_{\text{unsafe}}{=}2.0, z_{\text{safe}}{=}1.0 \Rightarrow
s{=}1.0 \Rightarrow$ unsafe-prob $\approx 0.73$). The score is \emph{not} a calibrated probability
until we calibrate it. \\
\textbf{Base vs.\ SFT / LoRA.} The \emph{base} is the instruction model before guard fine-tuning; SFT
(supervised fine-tuning) trains it on labeled safe/unsafe prompts; LoRA adds a small trainable adapter
instead of changing all weights. \emph{Base} means ``before the guard fine-tune,'' not ``untrained.''\\
\textbf{Average precision (AP), macro-AP.} A single number in $[0,1]$ for how well the score
\emph{ranks} truly-unsafe prompts above safe ones, across all cutoffs (area under the precision--recall
curve; no threshold needed). \emph{Macro}-AP averages AP across benchmarks with equal weight, so a big
benchmark can't dominate. (Why not accuracy? A 95\%-safe set scores 95\% by always saying ``safe.'')\\
\textbf{Represented-source vs.\ dataset-held-out (transfer).} \emph{Represented} = held-back rows from a
source the guard trained on; \emph{transfer} = datasets whose rows were never used in training. ``Held
out'' here means dataset-held-out, \emph{not} sealed-from-the-researcher.\\
\textbf{Operating point, FPR/TPR.} Picking one cutoff turns scores into decisions; TPR = fraction of
unsafe prompts caught, FPR = fraction of safe prompts wrongly flagged. \emph{Macro} averages each
benchmark's rate; \emph{pooled} lumps all rows first.\\
\textbf{Paired hierarchical bootstrap, interval, estimand.} To put a 95\% range on a change, we
re-draw the data many times (keeping each adapter paired with its own base, respecting the
checkpoint$\to$seed grouping) and re-compute. Intervals are conditional on this \emph{fixed panel}
(the exact quantity we estimate) --- they are not population claims or significance tests.\\
\textbf{Objectives.} SFT (cross-entropy on the correct verdict); \emph{DPO} (learn from a
correct-vs-wrong \emph{pair}); \emph{GRPO} (reinforcement learning with the label as a verifiable
reward --- no learned reward model). \\
\textbf{Composition.} Combining two guards by averaging their \emph{output} scores (not their internal
weights) --- portable but runs both models.
\end{background}

\begin{background}{Mortgage terms (Act IV), in plain words}
\textbf{HMDA:} the U.S.\ Home Mortgage Disclosure Act; lenders publicly report loan-level records
(purpose, amount band, action, denial reason). \textbf{Fair lending:} bars credit discrimination on
protected traits. \textbf{Redlining:} refusing/worsening credit for a whole neighborhood as a proxy for
race. \textbf{Disparate treatment} (intentional) vs.\ \textbf{disparate impact} (a neutral rule that
falls harder on a protected group). \textbf{Proxy/coded discrimination:} a stand-in (ZIP, language) for
a protected group. \textbf{ATR/QM:} ability-to-repay rule. \textbf{Adverse-action notice:} the required
true reason for a denial. \textbf{UDAAP:} unfair/deceptive/abusive practices. \textbf{Dual labels G,D:}
\textbf{G} = would an ordinary guard call this unsafe? \textbf{D} = would honoring it break mortgage
policy? \textbf{G0/D1} = looks safe, yet a violation --- the case that matters.
\end{background}

\section{The shared experimental setup}
\label{sec:setup}

Every study below reuses one fixed, purposively-chosen panel of four checkpoints and one frozen,
decontaminated 1{,}200-row training manifest, scored by the identical single-token
$z_{\text{unsafe}}-z_{\text{safe}}$ head and the same tie-aware macro-AP, with fail-closed provenance
locks binding data, code, and software versions. The \emph{same} four checkpoints recur across all four
acts; \textbf{Qwen3-4B is the recurring character} --- the strongest base, and (as we will see) the one
that specializes most, the one composition helps least, and the best-and-fairest zero-shot mortgage
guard. Reusing one panel and one manifest is what makes SFT, DPO, GRPO, composition, and the domain
evaluations directly comparable rather than three stapled results.

\section{Act I --- Fine-tuning specializes: large represented gain, no transfer}
\label{sec:actI}

\Cref{tab:primary} gives, per checkpoint, the untuned-base and mean-SFT macro-AP in each regime, the
paired base-to-SFT change, and the fixed-panel aggregate. Represented-source ranking rises sharply and
uniformly --- every SFT guard reaches $\approx 0.98$ regardless of where its base started --- for an
aggregate change of $\RepDelta$ [$\RepCILower,\RepCIUpper$]. Transfer barely moves on average
($\TransferDelta$ [$\TransferCILower,\TransferCIUpper$]) but the average \emph{hides a split}: the
change is negative for three checkpoints and positive for one, ranging from SmolLM2 ($+0.0400$) to
Qwen3-4B ($-0.1499$). \Cref{fig:plane} shows the signature directly: \SpecializationSeedCount{} of
\TotalSeedCount{} (checkpoint, seed) points fall in the lower-right ``specialize'' quadrant (represented
up, transfer down), \UniformGainSeedCount{} in uniform gain, none in uniform loss.

\begin{table}[H]\centering\small
\caption{Act I --- clean-v2 fixed-panel result: base and mean-SFT macro-AP by regime, paired
base-to-SFT deltas with two-sided 95\% hierarchical-bootstrap intervals, and the fixed-panel aggregate.
Descriptive under the locked precision-focused mode; conditional on this panel.}
\label{tab:primary}
\resizebox{\textwidth}{!}{\input{generated/tab_primary_gen}}
\end{table}

\begin{figure}[H]\centering
\includegraphics[width=0.6\textwidth]{specialization_plane.pdf}
\caption{Each point is one (checkpoint, seed): horizontal axis is the represented-source macro-AP
change, vertical axis the transfer change. \SpecializationSeedCount{}/\TotalSeedCount{} land in the
lower-right specialization quadrant --- the visual signature of Act I.}
\label{fig:plane}
\end{figure}

\begin{takeaway}{Why the $\approx\!-0.06$ transfer average is misleading}
It pools opposing per-checkpoint effects (SmolLM2 $+0.04$ vs.\ Qwen3-4B $-0.15$). The panel average is
small; the panel is \emph{not} uniform. And $\Delta = \text{SFT}-\text{base}$ is mathematically coupled
to base AP, so ``stronger bases specialize more'' is a hypothesis from four points, not a law.
Retrospective, estimation-only.
\end{takeaway}

\subsection{At a deployable threshold}
\label{sec:actI-op}
Ranking is threshold-free; deployment is not. At a calibration-selected operating point,
represented-source recall (TPR) rises from \RepBaseTPRPct\% to \RepSFTTPRPct\%, but on transfer the
false-alarm rate rises from \TransferBaseFPRPct\% (base) to \TransferSFTFPRPct\% (SFT), and single-class
stress shows HarmBench recall \emph{falling} from \HarmBenchBaseRecallPct\% to \HarmBenchSFTRecallPct\%
(\Cref{tab:sensitivity}). So the specialized guard raises more false alarms off-source and catches less
of a hard attack set --- the deployment cost behind the ranking story.

\begin{table}[H]\centering\small\setlength{\tabcolsep}{6pt}
\caption{Act I --- per-benchmark paired deltas (upper) and the calibration-targeted 5\% FPR operating
point plus single-class stress diagnostics (lower). Generated from the lock-bound scores.}
\label{tab:sensitivity}
\input{generated/tab_sensitivity_gen}
\end{table}

\section{Act II --- Does the objective matter? SFT vs DPO vs GRPO}
\label{sec:actII}

\begin{pending}{\textsc{pending} the locked Paper C run --- pre-registered, no numbers invented}
Act I used SFT. Does the \emph{choice of fine-tuning objective} change how much a guard specializes?
We run \textbf{DPO} and \textbf{GRPO} on the identical panel, manifest, seeds, and scorer (so the only
difference is the objective), and compare against the Act I SFT arm. The reward is the \emph{verifiable
label} --- no learned reward model. The hypotheses, metrics, and decision rules are \textbf{pre-registered}
(\code{docs/paper-c-prereg.md}, bound into the Paper C lock \emph{before} the run):
\textbf{H1} objectives differ mainly on transfer, ordered by how far each moves the model from base;
\textbf{the GRPO single-token null} --- because our action space is one token/two outcomes, the
exploration that lets RL generalize in the literature is absent, so we predict GRPO's transfer edge
over SFT is small or null (a bounded negative result is the expected, reportable outcome);
\textbf{H2} composition recovers transfer in proportion to how far the objective moved from base.
Results table and prose will replace this box once the locked run completes and its per-row scores are
committed; the analysis reuses the Act I macro-AP + bootstrap machinery unchanged.
\end{pending}

\section{Act III --- Compose, don't (re)tune: recovering transfer without retraining}
\label{sec:actIII}

Act I raises a question: rather than tuning harder, can we \emph{keep the base's judgment} in the
decision? The candidate rule is deliberately simple. For one input, run the base and one SFT adapter,
convert each raw score to a probability with a calibrator fit only on a development split, and report
the fixed equal-weight average:
\begin{equation}
s_{\mathrm{comp}}(x) = \tfrac12\, C_b\!\big(s_b(x)\big) + \tfrac12\, C_{a}\!\big(s_{a}(x)\big).
\label{eq:comp}
\end{equation}
Nothing is retrained; the $\tfrac12$ weights are fixed in advance (no tuning on test rows); it needs two
inference passes. This is \emph{output-space} composition --- unlike weight-space WiSE-FT / model soups
\citep{wortsman2022wiseft,wortsman2022modelsoups}, it needs only comparable scores, not interpolable
weights, but it runs both models. \Cref{tab:pilot-summary,tab:pilot-models} give the result: relative to
SFT, composition changes represented-source macro-AP by \PilotRepDeltaSFT{}
[\PilotRepDeltaSFTLow, \PilotRepDeltaSFTHigh] and transfer by \PilotTransferDeltaSFT{}
[\PilotTransferDeltaSFTLow, \PilotTransferDeltaSFTHigh] --- landing \emph{above the base} on transfer
(\PilotCompositionTransfer{} vs.\ \PilotBaseTransfer{}). Recovery relative to SFT is positive for all
four checkpoints, but the comparison with base is heterogeneous (it slightly \emph{hurts} the strongest
base, Qwen3-4B) --- so this is transfer \emph{recovery}, not Pareto dominance.

\input{generated/pilot_summary_table}
\input{generated/pilot_per_model_table}

\begin{takeaway}{Ranking $\neq$ calibration}
At a 5\%-FPR target, the composition's realized transfer false-alarm rate is
\PilotCompositionTransferFPRPct\% --- better than SFT (\PilotSFTTransferFPRPct\%) but still above target
(\Cref{tab:pilot-operating}). Recovering \emph{rank} does not deliver a transferable \emph{threshold};
a deployment claim from AP alone would miss this.
\end{takeaway}

\input{generated/pilot_operating_point_table}

\paragraph{What Act III does not establish.} Not Pareto dominance; not a proven mechanism; the pilot
lacks the equal-cost SFT+SFT control and actual WiSE-FT rescoring; the logit-average and convex-weight
variants were visible during development and are reported only as ablations, non-promotable. Whether
composition also recovers DPO/GRPO transfer (H2) is answered in \Cref{sec:actII} once that run lands.

\section{Act IV --- When the domain changes everything: four high-compliance domains}
\label{sec:actIV}

The strongest form of the thesis: in a regulated domain a general-safety benchmark can \emph{fail to see
the violation at all}, because the violating request reads as polite text. We study four
high-compliance domains with two complementary instruments.

\subsection{Mortgage --- a benchmark built in depth (dual-label G$\times$D)}
We built a fixed, HMDA-grounded benchmark whose unit is one incoming request carrying \emph{two}
independent labels: general-safety $G$ and mortgage-policy $D$. The load-bearing quadrant is
\textbf{G0/D1} --- reads safe, is a compliance violation. Requests are grounded in de-identified, banded
HMDA fact sheets (never verbatim records) via an agentic pipeline (\Cref{fig:pipeline}), and we add
protected-class \emph{minimal pairs} (identical except one protected token) as a fairness-invariance
gate. \Cref{tab:composition} gives the 994-row composition.

\begin{figure}[H]\centering
\includegraphics[width=0.52\textwidth]{pipeline.png}
\caption{The agentic mortgage-benchmark construction pipeline: real HMDA records $\to$ de-identified
banded fact sheets $\to$ plan/ground/generate/adversarially-mutate/rubric-bound-judge $\to$ accept-iff
label matches target $\to$ provenance + family-isolated split $\to$ the frozen release.}
\label{fig:pipeline}
\end{figure}

\input{generated/mortgage_composition_table}

Scored zero-shot (as shipped, prompt-only) the four base checkpoints rank mortgage-policy violations
only \emph{moderately} (\Cref{tab:baseline}, AP$\cdot$D $0.67$--$0.85$, at or above AP$\cdot$G): asking
for discrimination or fraud reads as ``unsafe'' to a general model even without a jailbreak, so $G$ and
$D$ are only \emph{partially} orthogonal, and none reaches AP$\cdot$D near $1$ --- much of the subtle
G0/D1 stratum is unresolved by ranking. The protected-pair gap $\Delta_{\mathrm{context}}$ is the
sharpest discriminator: Qwen3-4B is invariant ($\Delta_{\mathrm{context}}\!=\!0.000$, best and fairest)
while Qwen2.5-1.5B shows a large gap ($0.183$) --- a fairness signal AP alone misses. The fixed-threshold
operating point is deliberately \emph{not} tabulated: for these clustered-score guards it is knife-edge
(its G0/D1 catch count swung $>\!50$ rows across library versions), itself a finding that naive
threshold transfer is unreliable --- echoing Acts I and III.

% GENERATED by tools/emit_baseline_tex.py from out_eval/baseline_table.json -- do not edit.
\begin{table}[htbp]
\centering
\small
\caption{Baseline zero-shot instruction guards on the frozen benchmark (\texttt{public\_test}, 146 rows: 75 G0/D1, 6 G1, 3 protected pairs). AP is recomputed in the repo canonical environment from the committed per-row scores (exactly reproducible); the AP$\cdot$D column carries a 2{,}000-resample bootstrap 95\% CI. Threshold-free AP$\cdot$D is moderate (0.67--0.85) --- and its per-guard CIs overlap, so we do not rank guards by it: soliciting fraud/discrimination reads as unsafe even without a jailbreak, so G and D are only partially orthogonal. $\Delta_{\mathrm{context}}$ is the mean absolute protected-pair score gap (lower is more invariant) over only $n=3$ pairs, so its magnitude is uncertain and it is read as a direction, not a calibrated fairness estimate. AP$\cdot$final equals AP$\cdot$D on this frozen set because the G1/D0 cell is empty (every $G$-positive row is also $D$-positive), so the composed label reduces to $D$. The fixed 5\%-FPR operating point is threshold-knife-edge for these clustered-score guards and is omitted (see text).}
\label{tab:baseline}
\begin{tabular}{lrrrr}
\toprule
Guard & AP$\cdot$G & AP$\cdot$D (95\% CI) & AP$\cdot$final & $\Delta_{\mathrm{context}}$ \\
\midrule
qwen25\_15b\_base & 0.681 & 0.793\,[.71, .87] & 0.793 & 0.183 \\
qwen3\_4b\_base & 1.000 & 0.851\,[.78, .91] & 0.851 & 0.000 \\
smollm2\_17b\_base & 0.261 & 0.672\,[.57, .78] & 0.672 & 0.023 \\
smollm3\_3b\_base & 0.546 & 0.733\,[.64, .81] & 0.733 & 0.010 \\
\bottomrule
\end{tabular}
\\[2pt]{\footnotesize Not scored (gated, HF license not accepted): llama\_guard\_3\_1b, wildguard\_7b.}
\end{table}


\begin{takeaway}{Measuring stick, not a legal finding}
Mortgage labels are assigned by an LLM judge against written policy cards --- \emph{not}
SME-adjudicated (self-consistency, no human Fleiss-$\kappa$). The G1/D0 quadrant is empty and samples
are small. This benchmark \emph{surfaces} guard behavior; it certifies nothing about any real lender or
model.
\end{takeaway}

\subsection{Finance, health, law --- an external expert-annotated replication (ExpGuard)}
\begin{pending}{\textsc{pending} the ExpGuard base eval --- running; no numbers invented}
To test whether the specialization/transfer pattern recurs across regulated verticals \emph{with
expert labels} (a strictly stronger labeling tier than the mortgage LLM-judge), we score the same four
checkpoints on ExpGuard --- an expert-annotated moderation set spanning \textbf{finance}, \textbf{health},
and \textbf{law} (input-prompt classification, matching our task). We report aggregate and per-domain
AP/AUROC, base vs.\ tuned. This is single-label (not the mortgage dual G$\times$D construct) --- one
domain built in depth plus three external verticals for breadth. The per-domain table replaces this box
once the eval completes and its text-free per-row scores are committed.
\end{pending}

\section{Synthesis --- one lesson at four scales}
\label{sec:synthesis}
The four acts are one lesson. A guard's benchmark score is co-produced by the benchmark: SFT buys
represented-source ranking, not transfer (Act I); the specialization is a property to measure across
\emph{objectives}, not assume (Act II); keeping the base in an output-space average \emph{recovers} some
transfer for free of retraining, though not a threshold (Act III); and a change of domain can hide the
violation entirely, so regulated domains need domain-grounded evaluation and a fairness gate (Act IV).
The recurring cast makes it concrete: Qwen3-4B, the strongest base, specializes the most on transfer,
is the one composition hurts, yet is the best and fairest zero-shot mortgage guard --- the ranking
flips with the benchmark.

\section{The practitioner's decision guide}
\label{sec:practitioner}
\begin{enumerate}[leftmargin=1.6em,itemsep=1pt]
  \item \textbf{Always compare to your own base}, on represented \emph{and} held-out sets, plus a
  benign over-refusal and a hard-attack stress set --- a leaderboard delta vs.\ other models hides
  specialization.
  \item \textbf{Expect an objective to specialize} unless measured otherwise (Act II); don't assume RL
  ``generalizes'' on a single-token guard.
  \item \textbf{If transfer matters, try composition before tuning harder} --- but re-calibrate the
  threshold on the target regime (ranking recovery $\neq$ calibration transfer).
  \item \textbf{In a regulated domain, build domain-grounded, dual-labeled evaluation} and a
  protected-pair invariance check; a general-safety score does not cover compliance.
\end{enumerate}
Each step carries its caveat: these are estimates on a fixed panel, and the domain labels are a
measuring stick, not a verdict.

\section{How honest is this? Evidence ledger and what it does NOT establish}
\label{sec:honesty}
\begin{center}\small
\begin{tabular}{@{}p{0.20\textwidth}p{0.40\textwidth}p{0.34\textwidth}@{}}
\toprule
Result & Establishes & Key limits / would-upgrade \\
\midrule
Act I (SFT), III (composition), II (objective axis) & paired base-to-tuned change on a fixed panel &
retrospective, estimation-only, manifest inspected in dev; not confirmatory. Upgrade: prospective
uninspected cohort; SFT+SFT / WiSE-FT controls. \\
Mortgage benchmark & a dual-labeled, HMDA-grounded measuring stick for the G0/D1 stratum + fairness gate
& LLM-judge (not SME) labels; empty G1/D0; small single split. Upgrade: SME adjudication (Fleiss-$\kappa$). \\
ExpGuard (finance/health/law) & external, expert-annotated per-domain replication & single-label; a
transfer probe, not a dual-label construct. \\
\bottomrule
\end{tabular}
\end{center}
Two hard rules held throughout: retrospective and prospective numbers are never pooled, and no number is
upgraded to a causal, universal, or fair-lending claim.

\section{Validation roadmap}
\label{sec:roadmap}
(i) SME-adjudicate a stratified mortgage subset (report Fleiss-$\kappa$); (ii) lock a genuinely
uninspected prospective cohort for the objective axis; (iii) add the equal-cost SFT+SFT and WiSE-FT
composition controls; (iv) build dual-labeled finance/health benchmarks with expert sign-off; (v) score
gated open guards (Llama Guard 3, WildGuard) once licenses are accepted.

\section{Reproducibility}
\label{sec:repro}
Every number above is \code{\textbackslash input} from a committed generated table bound to a lock ---
none is hand-typed. A single entry point, \code{make reproduce} (\code{papers/unified-report/reproduce.py}),
regenerates all tables/figures from committed per-row scores by calling each study's analysis (Paper A in
its lock-pinned env; Paper B, mortgage, ExpGuard, and Paper C from committed scores), and asserts
byte-identity with the committed report tables. It needs no GPU and no network; only re-training the
guards or re-scoring the gated ExpGuard set requires a GPU / dataset access. Raw third-party rows are
referenced by pinned identifier + revision + content hash, not redistributed.

\section{Conclusion}
\label{sec:conclusion}
A small guard's benchmark score is not a fixed property of the guard --- the benchmark co-produces it.
Measured honestly, fine-tuning specializes rather than uniformly improves; the effect is worth measuring
across objectives; a retraining-free composition recovers some of the lost transfer; and in regulated
domains a general-safety score does not cover compliance. The disciplined next step is not a more
confident headline but a prospectively locked evaluation on genuinely uninspected data.

\bibliographystyle{plainnat}
\bibliography{refs}

\appendix
\section{Per-seed data and provenance}
\label{app:seeds}
\begin{table}[H]\centering\scriptsize\setlength{\tabcolsep}{2.5pt}
\caption{Observed represented-source and transfer macro-AP change for every SFT seed (Act I),
generated from the same validated score matrix as \Cref{tab:primary}.}
\label{tab:seed-values}
\input{generated/tab_seed_values_gen}
\end{table}

\end{document}
